\section{Formelzeichen und Abkürzungen}
\hspace{0.7cm}
\begin{tabularx}{15.1cm}{llX}
	\textbf{Formelzeichen} & \textbf{Einheit} & \textbf{Beschreibung} \\
	\hline
	$a_\text{e}$ & \unit{mm} & Eingriffsweite \bigstrut[t] \\ %bigstrud[t] und bigstrud[b] verwenden um Abstand zu linien korrekt abzubilden.
	$a_\text{p}$ & \unit{mm} & Schnitttiefe \\
	$f_\text{z}$ & \unit{mm} & Vorschub je Zahn \\
	$z$ & \unit{mm} & Zähnezahl  \bigstrut[b] \\[0.75cm]
	
	\textbf{Abkürzung} &  & \textbf{Beschreibung} \bigstrut[t] \\
	\hline
	DC &   & Gleichstrom \\
\end{tabularx}

%Sortiert wird alphabetisch, ungeachtet der Groß- und Kleinschreibung.
%Griechische Buchstaben werden nach den lateinischen angeführt.
%Dieses Beispiel zeigt wie man die Tabelle manuell aufbauen könnte.
%Es gibt jedoch auch Pakete die bei der Erstellung von Verzeichnissen unterstützen können.